\section{Discussion}

We asked whether building human cognitive or developmental constraints into a language model's architecture or training data brings it algorithmically closer to the human language system, or whether standard, unconstrained LLMs already achieve as much algorithmic convergence through task-scale optimization alone. Across seventeen datasets spanning five neural recording modalities and three behavioral measures, we find consistently for the latter. Standard LLMs match or exceed every constrained group on both BrainScore and BehavScore; the Pareto front trading off neural against behavioral predictive power is occupied exclusively by Standard LLMs, whether we look at all words or specifically at the syntactically hardest ones; and scale itself keeps paying off within the unconstrained group in a way it does not for the memory-constrained groups. This is the pattern the Platonic Representation Hypothesis predicts \citep{huh2024platonic}, and not the pattern the human-constraint hypothesis predicts: shared task demands, not shared inductive biases, appear to be what pulls these systems toward common representations.

The clearest result in support of that reading is architectural. Restricting a Transformer's attention to nearby context -- the intervention most directly motivated by a specific human processing theory, capacity-limited working memory \citep{lewisvasishth2005activation,futrell2020lossy} -- makes a model a \emph{worse} predictor of human neural and behavioral responses, at any training scale. If bounded working memory captured something about human processing that unconstrained models miss, we would expect the opposite. Restricting only the training data, while leaving the architecture untouched (BabyLM), fares better: competitive with Standard LLMs on BehavScore, and briefly ahead of them at the last layer specifically (Appendix~\ref{app:best-layer}), though never on NeuralScore. The two constraints do not compound -- Constrained-BabyLM is not better than either constraint alone -- and the constraint with an actual mechanistic motivation is the one that hurts most.

Two design choices shape how we read these results. We did not match model groups on parameter count or perplexity, because the question we ask is whether a constraint helps a model that is otherwise free to be as capable as it can, not whether it helps at fixed capability -- a smaller, worse-at-next-token-prediction memory-constrained model is exactly the comparison the working-memory hypothesis needs to win, and it does not. And we treat surprisal, not embeddings, as the primary comparison, because a scalar surprisal value carries no dimensionality advantage for larger models the way a wider hidden state does for ridge regression; that the Standard LLM advantage holds this cleanly on BehavScore is, if anything, the more decisive of our two results. We also note that our benchmark tests predictive power, not internal mechanism -- a necessary but not sufficient condition for algorithmic convergence.

Taken together, we read these results as pointing away from architectural constraints motivated by isolated psycholinguistic phenomena, and toward understanding what it is about task-scale optimization itself that produces human-like solutions as a byproduct. If bounded attention, the most theoretically motivated constraint we test, does not help, the more promising path may not be constraining a model's computation at all, but identifying which properties of standard training at scale are actually responsible for the convergence we observe.
